% !TeX encoding=utf-8
\documentclass[a4paper]{feidippp}
\usepackage[pdftex]{graphicx}
\DeclareGraphicsExtensions{.pdf,.png.,mps.,jpg}
\graphicspath{{figures/}}

\usepackage[utf8]{inputenc}
\usepackage[T1]{fontenc}
\usepackage[slovak]{babel}

\usepackage{lmodern}

\usepackage{amsmath,amssymb,amsfonts}

\def\figurename{Obrázok}
\def\tabname{Tabuľka}



%\usepackage[dvips]{graphicx}
%\DeclareGraphicsExtensions{.eps}



%\usepackage[pdftex]{hyperref}   %% tlac !!!
\usepackage[pdftex,colorlinks,citecolor=magenta,bookmarksnumbered,unicode,pdftoolbar=true,pdfmenubar=true,pdfwindowui=true,bookmarksopen=true]{hyperref}
\hypersetup{%
pdfcreator={pdfcsLaTeX},
pdfkeywords={Používateľská príručka, ACCESS-SECURITY},
pdftitle={Používateľská príručka --- Access-Security},
pdfauthor={Oleksandr Mykhailyshyn},
pdfsubject={Používateľská príručka}
}

% Citovanie podla mena autora a roku
%\usepackage[numbers]{natbib}
\usepackage{natbib} \citestyle{chicago}

%\usepackage{mathptm} %\usepackage{times}

\katedra{počítačov a informatiky}
\department{Department of Computers and Informatics}
\odbor{Informatika}
\autor{Oleksandr Mykhailyshyn}
\veduci{doc.~Ing.~Eva Chovancová, PhD.}
%\konzultant{}
\nazov{Zvýšenie bezpečnosti vybraných šifrovacích algoritmov v~prístupových systémoch}
\kratkynazov{Zvýšenie bezpečnosti vybraných šifrovacích algoritmov v~prístupových systémoch}
\nazovprogramu{ACCESS-SECURITY}
\klucoveslova{kontrola prístupu, RFID, AEAD, ChaCha20-Poly1305, HMAC, anti-replay, audit log}
\title{Enhancing Security of Selected Encryption Algorithms in Access Control Systems}
\keywords{access control, RFID, AEAD, ChaCha20-Poly1305, HMAC, anti-replay, audit log}
\datum{TODO.TODO.2026}



\begin{document}
\bibliographystyle{dcu}

\sloppy
\emergencystretch=2em
\tolerance=2000
\hyphenpenalty=500
\hbadness=10000



\titulnastrana

\newpage


\tableofcontents



\newpage

\setcounter{page}{1}

\section{Funkcia programu}

\emph{Access-Security} je prototyp systému kontroly fyzického prístupu na báze RFID/NFC. Hardvérová čítačka (ESP32 + MFRC522) prečíta UID karty a~odošle autentifikovanú HTTP požiadavku serveru, ktorý na základe roly karty a~oprávnení dverí rozhodne o~povolení alebo zamietnutí prístupu. Komunikácia je chránená HMAC podpisom, AEAD šifrovaním (ChaCha20-Poly1305), anti-replay mechanizmom a~detektorom protokolových anomálií; všetky udalosti sú zaznamenané v~tamper-evident auditnom logu.

Program je určený na laboratórne nasadenie a~experimentálnu validáciu kryptografických mechanizmov; produkčné nasadenie vyžaduje dodatočné kroky uvedené v~kapitole~\emph{Obmedzenia programu}. Detailný technický popis je v~systémovej príručke priloženej k~práci.

\paragraph{Autor:} Oleksandr Mykhailyshyn, FEI TUKE, 2026.

\paragraph{Licencia:} zdrojový kód aj všetky externé závislosti sú dostupné pod permisívnymi open-source licenciami (MIT, BSD, ISC, Apache~2.0).


\section{Súpis obsahu dodávky}

Dodávka pozostáva z~jedného elektronického média (USB / CD), ktoré obsahuje:

\begin{itemize}
\item zdrojový kód projektu (vetva \texttt{analytical} Git repozitára, vrátane produkčného serverového kódu, firmvéru ESP32, experimentálneho harnessu, analytického pipeline a~konfiguračných súborov);
\item \texttt{Dockerfile} pre kontajnerizované nasadenie;
\item dokumentáciu: text diplomovej práce, systémovú príručku a~túto používateľskú príručku (vo formáte PDF aj v~zdrojovom \LaTeX{} formáte);
\item súbor \texttt{README.md} v~koreňovom adresári so stručným prehľadom projektu.
\end{itemize}

\noindent Master kľúč (\texttt{secrets/master\_key.hex}) nie je súčasťou dodávky --- musí byť vygenerovaný pri prvom nasadení (popísané v~podkapitole~3.3).

\section{Inštalácia programu}

\subsection{Požiadavky na technické prostriedky}

Server: x86-64/ARM64, min.\ 2~jadrá, 2\,GB RAM, 1\,GB diskového priestoru pri preklade (150\,MB pre prevádzku). Hardvér: ESP32 DevKit, modul MFRC522, zelená LED, piezomenič a~USB kábel pre nahranie firmvéru.

\subsection{Požiadavky na programové prostriedky}

Linux Ubuntu 22.04+, GCC $\geq$~12, CMake $\geq$~3.20, \texttt{libsqlite3-dev}, Python $\geq$~3.10 (pre analytický pipeline). Pre firmvér Arduino IDE 2.x s~podporou ESP32 a~knižnicou MFRC522. Voliteľne Docker Engine 20.10+ pre kontajnerizované nasadenie.

\subsection{Vlastná inštalácia}

\paragraph{Server.} V~koreňovom adresári projektu:
\begin{verbatim}
mkdir build && cd build
cmake ..
cmake --build . -j$(nproc)
\end{verbatim}

\noindent Pri prvom zostavení CMake FetchContent automaticky stiahne externé závislosti (libsodium, cpp-httplib, yaml-cpp, nlohmann/json, Google Test). Verifikácia: \texttt{ctest --output-on-failure}.

\paragraph{Master kľúč.} Pred prvým spustením vygenerovať 32-bajtový kľúč:
\begin{verbatim}
mkdir -p secrets
openssl rand -hex 32 > secrets/master_key.hex
\end{verbatim}

\paragraph{Firmvér ESP32.} V~Arduino IDE upraviť \texttt{hw\_layer/hw\_config.h} (Wi-Fi SSID/heslo, URL servera, HMAC secret, identita čítačky), zvoliť board \emph{ESP32 Dev Module} a~nahrať skech \texttt{hw\_layer.ino} cez USB.

\paragraph{Docker (alternatíva).} \texttt{docker build -t access-security .} a~následne \texttt{docker run -p 8080:8080 -v ./data:/app/data -v ./secrets:/app/secrets access-security}.

\subsection{Popis štruktúry programu}

Po spustení servera (binárny súbor \texttt{access\_admin} s~argumentom YAML konfigurácie) sa otvoria dva HTTP porty: \texttt{8080} (frontend a~webové rozhranie) a~\texttt{8081} (DecisionService, len loopback). Webové administrátorské rozhranie je prístupné na adrese \texttt{localhost:8080/static/}\allowbreak\texttt{index.html}.

\paragraph{Typický tok spracovania požiadavky:}
\begin{enumerate}
\item Používateľ priloží kartu k~čítačke.
\item Firmvér prečíta UID, vypočíta HMAC podpis a~odošle HTTP požiadavku na port 8080.
\item Frontend služba overí HMAC, vytvorí AEAD frame a~odošle ho na DecisionService (port 8081).
\item DecisionService dešifruje frame, overí anti-replay, spustí anomálny detektor, vyhľadá kartu v~DB a~vráti rozhodnutie.
\item Frontend odošle JSON odpoveď firmvéru, ktorý signalizuje výsledok (zelená LED + krátke pípnutie pri povolení; dvojité pípnutie pri zamietnutí).
\end{enumerate}

\paragraph{Detailný popis architektúry, modulov a~tokov je v~systémovej príručke.}


\section{Použitie programu}

Po inštalácii a~spustení servera má používateľ dva spôsoby interakcie: koncový používateľ (držiteľ karty) komunikuje so systémom výhradne cez priloženie karty k~čítačke; administrátor systém spravuje cez webové rozhranie.

\subsection{Popis dialógu s~používateľom}

\paragraph{Koncový používateľ (držiteľ karty).} Po priložení karty čítačka signalizuje výsledok:
\begin{itemize}
\item \emph{Povolený prístup:} jedno krátke pípnutie (2\,kHz) so~svetlom zelenej LED.
\item \emph{Zamietnutý prístup:} dvojité pípnutie nižšieho tónu (500\,Hz).
\end{itemize}

\paragraph{Administrátor.} Otvorí v~prehliadači adresu \texttt{localhost:8080/static/}\allowbreak\texttt{index.html}, zadá administrátorský token (predvolene \texttt{admin}) a~prepína medzi sekciami pomocou navigačného menu.

\subsection{Popis funkcií a~podfunkcií programu}

Webové rozhranie ponúka šesť hlavných sekcií:

\begin{itemize}
\item \emph{Readers} --- registrácia čítačiek, priradenie verzie kľúča, zrušenie karantény.
\item \emph{Door roles} --- definícia rolí pre dvere (RBAC pravidlá).
\item \emph{Cards} --- pridávanie/odoberanie kariet a~ich rolí (UID sa ukladá ako HMAC odtlačok).
\item \emph{Audit} --- prehliadanie auditného logu a~kryptografická verifikácia integrity reťaze (auto-refresh každú sekundu).
\item \emph{Runtime logs} --- prevádzkové udalosti v~reálnom čase (polling 500\,ms).
\item \emph{DB Export/Import} --- záloha a~obnova databázy (s~kontrolou integrity \texttt{PRAGMA integrity\_check}).
\end{itemize}

\noindent Každá operácia sa vykonáva cez REST endpoint chránený hlavičkou \texttt{X-Admin-Token}; výsledok sa zobrazí v~príslušnej tabuľke alebo stavovom paneli.

\section{Popis vstupných, výstupných a~pracovných súborov}

\paragraph{Vstupné:}
\begin{itemize}
\item \texttt{config/access\_security.yaml} --- hlavná konfigurácia (port, anti-replay, decision mode).
\item \texttt{secrets/master\_key.hex} --- 32-bajtový master kľúč v~hex formáte (64 znakov).
\item \texttt{hw\_layer/hw\_config.h} --- konfigurácia firmvéru (kompiluje sa do binárneho súboru).
\end{itemize}

\paragraph{Výstupné a~pracovné:}
\begin{itemize}
\item \texttt{data/access.db} --- SQLite databáza (čítačky, karty, RBAC, audit log); vytvorí sa automaticky pri prvom spustení.
\item HTTP odpovede vo formáte JSON (\texttt{\{"allow": bool, "reason": string\}}).
\end{itemize}

\noindent Detailný popis formátov a~štruktúr je v~systémovej príručke.

\section{Obmedzenia programu}

Aktuálna verzia je laboratórny prototyp. Pred produkčným nasadením je potrebné brať do úvahy:

\begin{itemize}
\item \emph{Bez TLS:} komunikácia medzi čítačkou a~serverom prebieha cez nešifrovaný HTTP. HMAC zabezpečuje integritu a~autentifikáciu, ale nie dôvernosť.
\item \emph{Master kľúč na disku:} \texttt{master\_key.hex} je obyčajný súbor; pre produkciu sa odporúča HSM alebo vault.
\item \emph{Karta = identifikátor, nie kryptografický dôkaz:} klonovanie karty nie je riešené na úrovni protokolu.
\item \emph{Karanténa bez auto-zrušenia:} zakaranténovaná čítačka zostáva blokovaná až do manuálneho zásahu administrátora.
\item \emph{Sériový debug výstup ESP32:} firmvér vypisuje citlivé údaje (UID, HMAC podpis) na sériovú konzolu --- v~produkcii odporúčané vypnúť.
\item \emph{Bez offline režimu:} pri výpadku servera čítačka nemá lokálnu cache oprávnení (fail-closed).
\end{itemize}

\noindent Detailné zhodnotenie obmedzení a~smery ďalšieho vývoja sú v~systémovej príručke a~v~kapitole 4 hlavnej diplomovej práce.

\section{Chybové hlásenia}

Server pri zamietnutí požiadavky vracia v~JSON odpovedi pole \texttt{reason}, ktoré identifikuje príčinu. Najbežnejšie hodnoty:

\begin{itemize}
\item \texttt{ok} --- prístup povolený.
\item \texttt{unknown\_reader} / \texttt{unknown\_card} --- čítačka alebo karta nie je registrovaná.
\item \texttt{forbidden} --- karta nemá rolu povolenú pre dané dvere.
\item \texttt{replay} --- detekované opakovanie sekvenčného čísla.
\item \texttt{decrypt\_failed} --- AEAD overenie zlyhalo (poškodený frame, podvrhnutý nonce a~pod.).
\item \texttt{quarantined} --- čítačka je v~karanténe (zrušenie cez admin API).
\item \texttt{bad\_payload} / \texttt{bad\_action} --- frame má neplatný obsah.
\item \texttt{remote\_unavailable} --- DecisionService nedostupný (po~3 retry pokusoch, fail-closed).
\item \texttt{internal\_error} --- neočakávaná výnimka pri spracovaní (frame zamietnutý).
\end{itemize}

\noindent Chyby súvisiace s~operačným systémom (\texttt{listen failed}, chyby SQLite, sieťové výnimky) sa vypisujú na štandardný výstup servera.

\section{Príklad použitia}

Typická prvá konfigurácia:

\begin{enumerate}
\item Spustiť server cez binárny súbor \texttt{access\_admin} s~konfiguračným YAML súborom ako argumentom.
\item Otvoriť webové rozhranie na adrese \texttt{localhost:8080}, zadať admin token (\texttt{admin}).
\item V~sekcii \emph{Readers} pridať čítačku s~ID rovnakým ako \texttt{HW\_READER\_ID} v~\texttt{hw\_config.h}.
\item V~sekcii \emph{Door roles} pridať rolu (napr.\ \texttt{employee}) pre dvere s~ID rovnakým ako \texttt{HW\_DOOR\_ID}.
\item V~sekcii \emph{Cards} pridať kartu --- zadať UID v~hex formáte (uppercase) a~rolu \texttt{employee}.
\item Priložiť kartu k~čítačke: pri správnej konfigurácii sa rozsvieti zelená LED a~ozve sa krátke pípnutie.
\item Skontrolovať záznam v~sekcii \emph{Audit}.
\end{enumerate}

\noindent Pre opätovné spustenie experimentov: \texttt{./build/experiments/s1\_replay} (alebo iný scenár \texttt{s2}\,--\,\texttt{s7}); kompletný E2E beh cez \texttt{./experiments/run\_e2e.sh}.

\end{document}


